\documentclass[10pt]{article}
\usepackage[T1]{fontenc}
\usepackage[utf8]{inputenc}
\usepackage[french]{babel}
\usepackage{amsmath,amssymb,mathrsfs,mathpazo}
\usepackage{geometry}
\geometry{a4paper,margin=2cm}
\usepackage{graphicx,array,tabularx,tabularray,colortbl,multirow,multicol,booktabs,enumitem,pifont,ulem,fancybox,diagbox,cancel}
\usepackage{pstricks,pstricks-add,pst-plot,pst-tree,pst-node,pst-text,pst-func,pst-math,pst-bspline,pst-3dplot,pst-eucl,pst-solides3d}
\let\pscalculate\undefined
\usepackage{pst-bezier}
\usepackage{tikz,pgfplots,tkz-tab}
\usetikzlibrary{arrows,arrows.meta,patterns,calc,positioning,trees,decorations.markings,shapes.geometric,shapes.misc,shapes.arrows,intersections,decorations.pathmorphing,decorations.pathreplacing,backgrounds,shadows,babel}
\pgfplotsset{compat=1.18}
\usepgfplotslibrary{fillbetween}
\usepackage[np]{numprint}
\usepackage{listings,makecell,pas-tableur,twoopt}
\let\barre\undefined
\usepackage[pstricks]{tablvar}
\definecolor{darkgreen}{RGB}{0,100,0}
\providecommand{\red}{\color{red}}
\providecommand{\blue}{\color{blue}}
\providecommand{\green}{\color{green}}
\providecommand{\black}{\color{black}}
\providecommand{\white}{\color{white}}
\providecommand{\gray}{\color{gray}}
\providecommand{\cyan}{\color{cyan}}
\providecommand{\magenta}{\color{magenta}}
\providecommand{\yellow}{\color{yellow}}
\providecommand{\violet}{\color{violet}}
\providecommand{\orange}{\color{orange}}
\providecommand{\euro}{\texteuro}
\providecommand{\EUR}{\texteuro}
\providecommand{\ieme}{\textsuperscript{e}}
\providecommand{\degres}{\ensuremath{^\circ}}
\providecommand{\degree}{\ensuremath{^\circ}}
\providecommand{\vv}[1]{\overrightarrow{#1}}
\providecommand{\Coord}[1]{#1}
\providecommand{\decosix}{}
\providecommand{\starredbullet}{\ensuremath{\bullet}}
\pagestyle{empty}
\setlength{\parindent}{0pt}
\usepackage{adjustbox}\DeclareUnicodeCharacter{2236}{:}\DeclareUnicodeCharacter{2190}{\ensuremath{\leftarrow}}
\providecommand{\R}{\mathbb{R}}
\providecommand{\N}{\mathbb{N}}
\providecommand{\Z}{\mathbb{Z}}
\providecommand{\Q}{\mathbb{Q}}
\providecommand{\D}{\mathbb{D}}
\providecommand{\C}{\mathbb{C}}
\providecommand{\e}{{\mathrm{e}}}
\providecommand{\vect}[1]{\overrightarrow{#1}}
\providecommand{\vectt}[1]{\overrightarrow{\mathrm{#1}}}
\providecommand{\vvt}[1]{\overrightarrow{\mathrm{#1}}}
\providecommand{\np}[2][]{{#2~#1}}
\providecommand{\vv}[1]{{\overrightarrow{#1}}}
\renewcommand{\d}{{\mathrm{d}}}
\providecommand{\degres}{{^\circ}}
\providecommand{\degre}{{^\circ}}
\providecommand{\blue}{\color{blue}}
\providecommand{\red}{\color{red}}
\providecommand{\enskip}{\quad}
\providecommand{\llbracket}{\lbrack\!\lbrack}
\providecommand{\rrbracket}{\rbrack\!\rbrack}
\providecommand{\euro}{{\text{€}}}
\providecommand{\no}{{\text{n°}}}
\providecommand{\interval}[2]{{[#1;#2]}}
\providecommand{\intervFO}[2]{{[#1;#2[}}
\providecommand{\dotfill}{\ldots}

\providecommand{\ds}{\displaystyle}
\providecommand{\textding}[1]{#1}
\providecommand{\Oij}{\(\mathrm{O};\vec{\imath},\vec{\jmath}\)}
\providecommand{\Oijk}{\(\mathrm{O};\vec{\imath},\vec{\jmath},\vec{k}\)}
\providecommand{\Ouv}{\(\mathrm{O};\vec{u},\vec{v}\)}
\providecommand{\intervFF}[2]{\ensuremath{[#1;#2]}}
\renewcommand{\euro}{\ifmmode\text{\texteuro}\else\texteuro\fi}
\DeclareUnicodeCharacter{00B0}{\ifmmode^{\circ}\else\textdegree\fi}
\DeclareUnicodeCharacter{20AC}{\ifmmode\text{\texteuro}\else\texteuro\fi}
% Énoncés : Éducation nationale. Transcription et corrigés : APMEP, Denis Vergès et contributeurs. Découpage et classement : CoopMaths, Sébastien Lozano, Éric Elter et contributeurs. Adaptation de la présentation : Maths Explorer.
% CC BY-SA : https://coopmaths.fr/www/about/
% https://coopmaths.fr/alea/static/bac/2026/tex/bac_2026_06_sujet2_antilles_2_cor.tex
\begin{document}
\noindent\textit{APMEP / CoopMaths. CC BY-SA. Présentation adaptée par Maths Explorer.}\par\medskip
\subsection*{Partie A}
\begin{enumerate}
\item 
	\begin{enumerate}
			\item 
			
			
			Comme le repère est orthonormé :
			\begin{itemize}
				\item $\mathrm{AC} = \sqrt{\left(\dfrac{9}{4}-0\right)^2 + (3-0)^2 + (0 - 0)^2} =\sqrt{\dfrac{81}{16}+ 9} =\sqrt{ \dfrac{81}{16} + \dfrac{144 }{16}} = \sqrt{\dfrac{225}{16}}=\dfrac{15}{4} $;
				\item $\mathrm{CB} = \sqrt{\left(\dfrac{9}{4}-\dfrac{9}{2}\right)^2 + (3-0)^2 + (0 - 0)^2} =\sqrt{\left(-\dfrac{9}{4}\right)^2+ 9} =\sqrt{ \dfrac{81}{16} + \dfrac{144 }{16}} = \sqrt{\dfrac{225}{16}}=\dfrac{15}{4} $;
			\end{itemize}
			
			Les longueurs sont égales, donc le triangle est bien isocèle en C.
			
			\item  \begin{itemize}
				\item  Dans un triangle isocèle, la hauteur issue du sommet principal est aussi une médiane (et une médiatrice).
			
			La hauteur issue de C est donc le segment [CI], où I est le milieu de [AB].
			
			I a donc pour coordonnées :\quad I$\left(\dfrac{x_{\mathrm{A}} + x_{\mathrm{B}}}{2} ; \dfrac{y_{\mathrm{A}} + y_{\mathrm{B}}}{2} ; \dfrac{z_{\mathrm{A}} + z_{\mathrm{B}}}{2}\right) $\quad soit \quad $\left(\dfrac{9}{4} ; 0 ; 0\right)$.
			
			\item L'aire d'un triangle est donnée par :\quad $\mathcal{A} = \dfrac{\text{base} \times \text{hauteur}
			}{2}$.
			
			Ici, on choisit comme base le segment [AB] et la hauteur correspondante, [CI].
			
			\item La distance AB est :\quad $\mathrm{AB} = \sqrt{\left(\dfrac{9}{2}-0\right)^2 + (0-0)^2 + (0 - 0)^2} =\dfrac{9}{2} $;
			
			La distance CI est :\quad $\mathrm{CI} = \sqrt{\left(\dfrac{9}{4}-\dfrac{9}{4}\right)^2 + (3-0)^2 + (0 - 0)^2} =\sqrt{ 9} =3$;
			\end{itemize}
			On a donc : \quad $\mathcal{A}_{\mathrm{ABC}} = \dfrac{\mathrm{AB} \times \mathrm{CI}}{2} = \dfrac{1}{2}  \times \dfrac{9}{2} \times 3 = \dfrac{27}{4} = 6,75$.
			
			
			
			L'aire du triangle  ABC est bien égale à $\np[m^2]{6,75}$.
		\end{enumerate}

	\item \begin{enumerate}
		\item On a : $\vectt{AB} \begin{pmatrix} \dfrac{9}{2}\\0\\0\end{pmatrix} $.
		et  \quad $\vectt{AC} \begin{pmatrix}\dfrac{9}{4} \\ 3\\0 \end{pmatrix} $.
		
		Ces vecteurs ont des coordonnées non proportionnelles, donc ils ne sont pas colinéaires.
		
		Les points A, B et C ne sont pas alignés, donc ils définissent un plan.
		
		De plus $z_{\mathrm{A}}=z_{\mathrm{B}}=z_{\mathrm{C}}=0$ 
		
		
	donc une équation cartésienne du plan (ABC) est bien :\quad $z=0$.
		
		\item  H est le projeté orthogonal du point D sur le plan (ABC) si et seulement si 
		
		$\begin{cases}
			H \in (\mathrm{ABC})\\
			\text{ la droite (HD) est orthogonale au plan (ABC)} \end{cases}$. 
			
			\begin{itemize}
				\item $z_{\mathrm{H}}=0$  donc $H \in (\mathrm{ABC})$
				
				\item  	$\vectt{HD} \begin{pmatrix}
					x_{\mathrm{D}} - x_{\mathrm{H}}\\
					y_{\mathrm{D}} - y_{\mathrm{H}}\\
					z_{\mathrm{D}} - z_{\mathrm{H}} \end{pmatrix} $
				\quad soit \quad $\vectt{HD} \begin{pmatrix} \dfrac{3}{2}-\dfrac{3}{2}\\1-1\\3-0 \end{pmatrix} $
				\quad finalement :\quad $\vectt{HD} \begin{pmatrix} 0\\0\\3\end{pmatrix} $.
				
				Or, une équation cartésienne du plan (ABC) étant $z=0$, un vecteur normal au plan (ABC) est  $\vect{m}\begin{pmatrix} 0\\0\\1 \end{pmatrix}$
				
				Les vecteurs $\vect{m}$ et $\vectt{HD}$ sont colinéaires, donc la droite (HD) est orthogonale au plan (ABC).
				
			\end{itemize} 
		
		
		Donc le point H$\left(\dfrac{3}{2} ; 1 ; 0\right)$ est le projeté orthogonal du point D sur le plan (ABC).
	\end{enumerate}
	\item  Avec les questions précédentes, on a calculé l'aire du triangle ABC, qui est l'une des quatre bases possibles pour le tétraèdre.
	
	Si on choisit ABC comme base de ABCD, alors la hauteur correspondante est la distance séparant le quatrième sommet (ici D), du plan contenant la base ABC, c'est-à-dire le plan (ABC).
	
	Cette distance est donc la distance séparant D de son projeté orthogonal sur (ABC), c'est-à-dire H, dont on a déterminé les coordonnées à la question \textbf{2 b.}.
	
	$\mathrm{DH} = \sqrt{0^2 + 0^2 + 3 ^2 } = 3$.
	
	On a donc : \quad $\mathcal{V}_{\mathrm{ABCD}} = \dfrac{1}{3} \times \mathcal{A}_{\mathrm{ABC}} \times \mathrm{DH} = \dfrac{1}{3} \times 6,75 \times 3 = 6,75$.
	
	Le volume du module est égal  à  $\np[m^3]{6,75}$.

\end{enumerate}


\subsection*{Partie B}
\begin{enumerate}
 \item 
	\begin{enumerate}
		\item 
		Déterminons les coordonnées des vecteurs :
		
		$\vectt{BC} \begin{pmatrix}
			x_{\mathrm{C}} - x_{\mathrm{B}}\\
			y_{\mathrm{C}} - y_{\mathrm{B}}\\
			z_{\mathrm{C}} - z_{\mathrm{B}} \end{pmatrix} $
		\quad soit \quad $\vectt{BC} \begin{pmatrix}\dfrac{9}{4} - \dfrac{9}{2} \\3-0\\0-0 \end{pmatrix} $
		\quad finalement :\quad $\vectt{BC} \begin{pmatrix} -\dfrac{9}{4}\\3\\0\end{pmatrix} $.
		\quad De même, \quad $\vectt{BD} \begin{pmatrix}  -3\\1\\3 \end{pmatrix} $.
		
		Ces vecteurs ont des coordonnées non proportionnelles, donc ils ne sont pas colinéaires.
		
		Les points B, C et D ne sont pas alignés, donc ils définissent un plan.
		
		On est dans un repère orthonormé, donc on peut utiliser la formule analytique du produit scalaire :
		
		\begin{itemize}
			\item $\vect{n} \cdot \vectt{BC} = 4 \times \dfrac{-9}{4} + 3 \times  3 + 3  \times  0 = -9 + 9 +0 = 0$;
			\item $\vect{n} \cdot \vectt{BD} = 4 \times (-3) +3 \times  1 + 3  \times  3 = -12 + 3 + 9 = 0$;
			\item $\vectt{BC}$ et $\vectt{BD}$ sont deux vecteurs non colinéaires qui forment une base du plan (BCD).
				\end{itemize}
			$\vect{n}$ est orthogonal à ces deux vecteurs, donc c'est un vecteur normal du plan (BCD).
			
		
		\item Si $\vect{n}\begin{pmatrix} 4\\3\\3 \end{pmatrix}$ est un vecteur normal au plan  (BCD), alors, par propriété, une équation du plan  (BCD) peut avoir la forme :\quad $4\times x + 3 \times y + 3 \times z + d = 0$, \quad où $d$ est un réel.
		
		Comme B est un point de (BCD), on doit avoir :\quad $4x_{\mathrm{B}} +3 y_{\mathrm{B}} + 3z_{\mathrm{B}} + d =0$.
		
		$\begin{aligned}
			4x_{\mathrm{B}} +3 y_{\mathrm{b}} + 3z_{\mathrm{B}} + d =0
			&\iff 4 \times \dfrac{9}{2} +3 \times 0 +3\times 0 +  d= 0\\
			&\iff 18 + d = 0\\
			&\iff d = -18
		\end{aligned}		$
		
		Une équation cartésienne de (BCD) est donc bien :\quad $4x+3y+3z - 18 = 0$.

		
		\item 
		Si la droite $\Delta$ est orthogonale au plan (BCD), alors, par définition, elle doit être dirigée par un vecteur normal au plan (BCD), donc ici, on en déduit que le vecteur $\vect{n} \begin{pmatrix} 4\\3\\3 \end{pmatrix}$ dirige la droite $\Delta$.
		
		Comme la droite $\Delta$ passe par le point D$\left(\dfrac{3}{2} ; 1 ; 3\right)$, une représentation paramétrique de la droite $\Delta$ est donnée par :\quad $\begin{cases}
			x=x_{\mathrm{D}} + x_{\vect{n}}\times t\\
			y=y_{\mathrm{D}} + y_{\vect{n}}\times t\\
			z=z_{\mathrm{D}} + z_{\vect{n}}\times t
		\end{cases}\text{avec } t \in \R$
		
		En remplaçant par les valeurs connues, on a :\quad $\begin{cases}
			x=\dfrac{3}{2} +4 t\\
			y=1 +3t\\
			z=3+3t
		\end{cases}\text{avec } t \in \R$
		
		
		
	\end{enumerate}
	
	
	\item Le point L  est l'intersection de la droite $\Delta$ et du plan (ABC).
	
	Comme c'est un point de la droite $\Delta$, il existe un réel $t_0$ tel que les coordonnées de ce point sont $\left( \dfrac{3}{2} + 4t_0 ; 1+3t_0 ; 3+3t_0 \right) $.
	
	C'est aussi un point de (ABC) dont ses coordonnées vérifient l'équation du plan.
	
	Autrement dit, $t_0$ est solution de l'équation :
	
	$\begin{aligned}[t]
		 3 + 3t_0  = 0
		&\iff 3t_0 = -3 \\
		&\iff t_0 = -1
	\end{aligned}$
	
	Le point L a donc pour coordonnées :\quad $\left( \dfrac{3}{2} - 4  ; 1-3  ; 3-3  \right) $ \quad  soit $\mathrm{L}\left(-\dfrac{5}{2} ; -2 ; 0\right) $.

	
	\item  Afin de déterminer l'angle $\widehat{\mathrm{LDH}}$, calculons le produit scalaire $\vectt{DL} \cdot \vectt{DH}$ de deux manières différentes. 
	
On a  $\vectt{DH} \begin{pmatrix} 0 \\ 0\\ -3\end{pmatrix} $.
	\quad et  \quad $\vectt{DL} \begin{pmatrix} -4  \\ -3\\ -3 \end{pmatrix} $.
	
	
	
	Nous sommes dans un repère orthonormé, nous pouvons utiliser la formule analytique : 
	
	$\begin{aligned}  \vectt{DH} \cdot \vectt{DL} &= x_{\mathrm{DH}}\times x_{\mathrm{DL}} + y_{\mathrm{DH}}\times y_{\mathrm{DL}} +z_{\mathrm{DH}}\times z_{\mathrm{DL}}\\
		&= 0\times (-4) + 0 \times (-3) -3 \times (-3)\\
		&= 9
		\end{aligned}$. 
		
		
		
		
		On a aussi : 
$	\vectt{DH} \cdot \vectt{DL}=\left\|\vectt{DH}\right\|\times\left\|\vectt{DL}\right\|\times \cos(\widehat{\mathrm{LDH}})$

Calculons la longueur  DL : $\mathrm{DL} = \sqrt{ (-4)^2+3^2+3^2} =\sqrt{ 16+9+9} = \sqrt{34} $ 

En remplaçant dans la formule précédente on obtient : 


$9 =3 \times \sqrt{34} \times \cos(\widehat{\mathrm{LDH}})$

D'où : $ \cos(\widehat{\mathrm{LDH}}) = \dfrac{9} {3\sqrt{34}}= \dfrac{3 \sqrt{34}}{34}$


donc $\widehat{\mathrm{LDH}} \approx  59,04$\,\textdegree{}.

 L'angle $\widehat{\mathrm{LDH}}$ est strictement inférieur à $60$\,\textdegree{} donc le module est bien  adapté aux débutants. 
\end{enumerate}

\end{document}
